\documentclass[12pt]{article}
\usepackage[utf8]{inputenc}
\usepackage[T1]{fontenc}
\usepackage{hyperref}
\title{Dampelas Language Resources for K-12 Curriculum Development}
\author{Compiled by Claude Sonnet 5.5 (claude-sonnet-5-5) for Andr\'as Kornai}
\date{}
\begin{document}
\maketitle

\section{Introduction}
Dampelas (ISO 639-3 \texttt{dms}) is an Austronesian language of Central Sulawesi, Indonesia, spoken on the north peninsula facing the Makassar Strait. Joshua Project, citing Ethnologue 2016, gives about 13,000 speakers, well above the 2,000 in the reference description, so the figure is uncertain. The search found almost nothing about the language itself.

\section{Primary Educational Resources}
No primers, textbooks or curricula in Dampelas were found.

\section{Linguistic and Reference Materials}
No grammar, dictionary or orthography was found in the search results.

\section{Religious and Community Resources}
\begin{itemize}
\item Joshua Project profile of the Dampelas people (reports written status unknown and a predominantly Muslim community): \url{https://joshuaproject.net/people_groups/11484}
\end{itemize}

\section{Digital and Audio Resources}
No audio, video or software resources were found.

\section{Curriculum Development Context}
Dampelas is also the name of a district in Donggala Regency (formerly Damsol): \url{https://en.wikipedia.org/wiki/Donggala_Regency}. Local primary schools exist there, but nothing indicates that they teach the language. A curriculum would need primary documentation work and community partnership.
\end{document}